\section{Related Work}

We organize prior work into four families directly relevant to the claims in this paper: the N-CMAPSS dataset and its baseline competitions, attention/transformer architectures for RUL, multi-scale temporal-convolution methods, and asymmetric losses for safety-critical prognostics. We close by situating our contribution against this literature, noting up front that nearly all of the attention- and multi-scale architectures we benchmark have, to date, been published only on the older C-MAPSS dataset and that head-to-head comparability with our N-CMAPSS results is therefore limited.

\subsection{N-CMAPSS dataset and baseline competitions}

The N-CMAPSS dataset~\cite{ariaschao2021ncmapss} extends the original C-MAPSS turbofan simulator~\cite{saxena2008damage} to run-to-failure trajectories generated under recorded flight conditions, producing variable-length flights on the order of $10^{4}$--$10^{5}$ timesteps per unit. The 2021 PHM Society Data Challenge introduced the first comparable leaderboard on this dataset. The challenge winner~\cite{lovberg2021variable} proposed a dilated-convolutional network explicitly designed to ingest variable-length input sequences with a learned condition-normalization step, which is the closest prior art to the variable-length flight handling we adopt. The second-place entry~\cite{devol2021inception} used inception modules to capture multiple temporal receptive fields from windowed N-CMAPSS data, and the third-place entry~\cite{solismartin2021stacked} reported a stacked DCNN trained on $L_w = 161$ windows, achieving validation L2 RMSE 6.24 and test NASA score 3.651. Outside the challenge, a hybrid physics-plus-deep-learning baseline~\cite{ariaschao2022fusing} reports an extended prediction horizon over a purely data-driven model on a fleet of nine N-CMAPSS engines, and the ANN-Flux baseline of Cohen et al.~\cite{ncmapss_eventual2023} reports RMSE 7.75 / NASA score 4.34 on N-CMAPSS. These are the primary numerical anchors for ``what is a good N-CMAPSS result,'' and they collectively establish a precedent that compact windowed inputs are the operating regime of choice on this dataset.

\subsection{Attention and transformer architectures for RUL}

The transformer~\cite{vaswani2017attention} replaces recurrence with self-attention as the primary sequence-modeling primitive, and Informer~\cite{zhou2021informer} adapts it to long sequences via ProbSparse attention and self-attention distillation, motivating input compression rather than raw long inputs. Two recent attention models are directly relevant. The STAR transformer~\cite{fan2024star} is a two-stage hierarchical encoder--decoder that applies sequential sensor-wise and temporal attention; it operates on $32$--$64$-step inputs and reports RMSE 10.61 / 13.47 / 10.71 / 15.87 on C-MAPSS FD001--FD004. TTSNet~\cite{ttsnet2025} fuses a transformer branch, a TCN branch and a multi-head self-attention branch over a noise-smoothed input and reports RMSE 11.02 / 13.25 / 11.06 / 18.26 on the same C-MAPSS subsets. Both are evaluated only on C-MAPSS, so their RMSE numbers are not directly comparable to ours: N-CMAPSS sequences are one to two orders of magnitude longer per unit, and the resulting RUL distributions differ.

\subsection{TCN and multi-scale temporal methods}

The temporal convolutional network~\cite{bai2018tcn} establishes dilated causal convolutions as a competitive sequence-modeling alternative to recurrent networks, with longer effective memory and faster training. The closest prior work to the MSTCN architecture evaluated here is the attention-based multi-scale TCN of Xu, Zhang and Miao~\cite{xu2024msattn}, which combines a multi-scale TCN backbone with self-attention at the head and a global-fusion-attention block at the tail. Their architecture is essentially the same family represented by our MSTCN implementation and is therefore the strongest direct precedent for our multi-scale temporal-attention comparison. As with the transformer literature, however, \cite{xu2024msattn} reports results on C-MAPSS rather than N-CMAPSS, so the architectural lineage is direct but the numerical comparison with our results is not.

\subsection{Loss functions for safety-critical RUL}

The dominant evaluation metric for RUL on C-MAPSS and N-CMAPSS is the asymmetric exponential PHM score introduced by Saxena et al.~\cite{saxena2008damage}, which penalizes late predictions exponentially harder than early ones (canonical penalty constants 13 and 10). Training-time loss functions that mirror this asymmetry have been proposed in recent work; the multiple-penalty-mechanism loss of \cite{mpm2024loss} pairs RMSE with similarity-based late-prediction penalties and reports improved advanced-prediction probability on C-MAPSS subsets. We adopt asymmetric MSE with a $2\times$ late-prediction penalty as a methodological choice anchored on \cite{saxena2008damage} and consistent with \cite{mpm2024loss}, rather than as an empirical contribution.

\subsection{Where this work fits}

Two observations motivate our contribution. First, recent attention and multi-scale-TCN architectures for RUL~\cite{fan2024star,ttsnet2025,xu2024msattn} are reported almost exclusively on the older C-MAPSS dataset, even though N-CMAPSS~\cite{ariaschao2021ncmapss} has been available since 2021 and exposes substantially longer, multi-condition, variable-length flights. Second, the published N-CMAPSS literature itself~\cite{lovberg2021variable,devol2021inception,solismartin2021stacked,ncmapss_eventual2023,ariaschao2022fusing} consists of a small number of competition entries and hybrid baselines, with no broad architecture sweep under matched preprocessing and training. Our work therefore benchmarks attention, multi-scale-TCN, transformer, WaveNet and RNN baselines on N-CMAPSS itself; the limited apples-to-apples N-CMAPSS attention literature is itself part of what makes the sweep informative. Where we cite C-MAPSS-only RMSE numbers (e.g.\ \cite{fan2024star,ttsnet2025,xu2024msattn}), we do so as architectural precedent, not as a direct numerical baseline.
